\chapter*{Úvod}
\addcontentsline{toc}{chapter}{Úvod}
\markboth{Úvod}{Úvod}

Za posledných päť rokov sme mohli byť svedkami radikálneho technologického posunu a transformácie vo svete informačných technológií.
Či už ide o nevinný veľký jazykový model (LLM) GPT-3 z roku 2020, ktorý odštartoval ďalší \enquote{boom} vo výskumnej oblasti umelej
inteligencie, kedy v roku 2026 môžu amatéri použiť \enquote{agentic} zručnosti najnovších LLM na podniknutie kybernetických útokov;
alebo v dôsledku globálnej pandémie vzostup práce z domu (\emph{homeoffice}), technológii VPN a digitálnych služieb; alebo mnohé
údajne kybernetické útoky\footnote{Ktoré sú súčasťou hybridnej vojny.} súvisiace s Rusko-­ukrajinskou vojnou, cielené na štátne weby
a systémy rôznych krajín sveta.

Aj vo svetle týchto udalosti sa kybernetická a informačná bezpečnosť (KIB) stáva omnoho väčšou nevyhnuteľnosťou pre organizácie a štátne
inštitúcie pôsobiace v digitálnom priestore. Bezpečnosť informačných systémov už nestačí riešiť iba na úrovni kódu a vývoja
bezpečného softvéru, ale vzrastá potreba riadenia prístupu a segmentácie systémov a sieti, aby sme minimalizovali dopady hrozieb.
Aby sme mohli riešiť KIB účinným a efektívnym spôsobom, potrebujeme integrovať KIB do všetkých úrovni organizácie systematicky pomocou
systému riadenia informačnej bezpečnosti (ISMS), ktorý vynucujeme politikou KIB. V organizácií potrebujeme zaviesť elementárny poriadok,
zdokumentovať všetky (kritické) aktíva, stanoviť bezpečnostné požiadavky vzhľadom na základné hodnoty CIA, identifikovať relevantné
hrozby, vyhodnotiť potenciálne riziká a zaviesť bezpečnostné opatrenia, aby sme zabezpečili želanú úroveň bezpečnosti v organizácií.

V tejto diplomovej prácu sa venuje analýze rizík, ktorá je kritickou súčasťou ISMS. Vychádzajúc z nemeckého IT-Grdunschutz-u a
príručke o analýze rizík \cite{BSI3}, popíšeme potrebný kontext, vysvetlíme celý proces, navrhneme a implementujeme \enquote{proof-of-concept}
systém na podporu analýzy rizík. Touto prácou pokračujeme a obsahovo nadväzujeme na našu predchádzajúcu bakalársku prácu \cite{BP},
kde sme analyzovali zákony \cite{ZoKB}, \cite{ZoITVS} a súvisiace vyhlášky\footnote{Analyzovali sme legislatívne povinnosti, ktoré
vyplývajú pre prevádzkovateľov základných služieb a poskytovateľov ISVS.}, popísali sme ISMS (podľa \cite{BSI1} a \cite{BSI2}) a
navrhli sme jednoduchý systém, ktorý má pomôcť zaviesť ISMS pre malé ISVS. 